\chapter{Integrale curbilinii}
\index{integrale!curbilinii}

\section{Elemente de teorie}

\textbf{Integrale curbilinii de speța întîi}

Fie $ \gamma = \gamma(t) $ o curbă netedă, definită pe un interval 
$ t \in [a, b] $. Se definește integrala curbilinie a unei
funcții $ f : \RR^3 \to \RR, f = f(x, y, z) $ prin formula:
\[
  \int_\gamma f(x, y, z) ds = \int_a^b f(x(t), y(t), z(t)) \cdot %
  \sqrt{x'(t)^2 + y'(t)^2 + z'(t)^2} dt.
\]

De asemenea, cîteva cazuri particulare de interes sînt:
\begin{itemize}
\item \emph{lungimea curbei $ \gamma $} se obține pentru $ f = 1 $,
  deci:
  \[
    \ell(\gamma) = \int_a^b \sqrt{x'(t)^2 + y'(t)^2 + z'(t)^2} dt;
  \]
\item dacă funcția $ f $ reprezintă \emph{densitatea} unui fir pe care
  îl aproximăm cu o curbă netedă $ \gamma $, atunci \emph{masa firului}
  se calculează cu formula:
  \[
    M(\gamma) = \int_\gamma f(x,y,z) ds;
  \]
\item în aceeași ipoteză de mai sus, \emph{coordonatele centrului de %
    greutate} al firului, $ x_i^G $ se calculează cu formula
  (am notat $ x_1 = x, x_2 = y, x_3 = z $):
  \[
    x_i^G = \dfrac{1}{M} \int_\gamma x_i f(x_1, x_2, x_3) ds,
  \]
  unde $ M $ este masa calculată mai sus.
\end{itemize}

\vspace{0.5cm}

\textbf{Integrala curbilinie de speța a doua}

Fie $ \omega = Pdx + Qdy + Rdz $ o 1-formă diferențială. Se definește
integrala curbilinie a formei $ \omega $ în lungul curbei $ \gamma $ ca mai sus
prin formula:
\[
  \int_\gamma \omega = \int_a^b (P \circ \gamma) x' + (Q \circ \gamma) y' + %
  (R \circ \gamma) z' dt,
\]
unde $ \gamma = \gamma(t), t \in [a, b] $ este o parametrizare a curbei $ \gamma $.

O aplicație fizică importantă a integralelor curbilinii de speța a doua
este calculul \emph{circulației} cîmpurilor vectoriale.

Fie, deci, un cîmp vectorial în spațiu:
\[
    \vec{V} = P(x, y, z)\vec{i} + Q(x, y, z)\vec{j} + R(x, y, z) \vec{k}.
\]
Atunci \emph{circulația\footnotemark} cîmpului vectorial în lungul unei curbe netede $ \gamma $
se calculează cu integrala curbilinie:
\[
    \int_\gamma P(x, y, z) dx + Q(x, y, z) dy + R(x, y, z) dz
\]
și putem asocia o 1-formă diferențială cîmpului, anume:
\[
    \omega = P(x, y, z) dx + Q(x, y, z)dy + R(x, y, z) dz.
\]

\footnotetext{Circulația unui cîmp vectorial este analogul fluxului, dar
    în dimensiune 1. Astfel, dacă fluxul se calculează pentru cîmpuri care
    traversează o suprafață (vom întîlni conceptul în lecția despre integrale
    de suprafață), circulația se calculează pentru cîmpuri în lungul
unor curbe.}

%%%%%%%%%%%%%%%%%%%%%%%%%%%%%%%%%%%%%%%%%%%%%%%%%%%%%%%%%%%%%%%%%%%%%%
\section{Formula Green-Riemann}
\index{formula!Green-Riemann}

Această formulă, care face parte din formulele integrale esențiale pe care
le vom studia, ne permite să transformăm o integrală curbilinie de speța a doua
într-una dublă.

Mai precis, avem formula:
\[
  \int_\gamma Pdx + Qdy = %
  \iint_K \frac{\partial Q}{\partial x} - \frac{\partial P}{\partial y} dxdy,
\]
unde $ K $ este suprafața bidimensională închisă de curba $ \gamma $.
\index{teorema!Green-Riemann}

\textbf{Observație importantă:} Pentru a putea aplica formula Green-Riemann,
este necesar ca drumul $ \gamma $ să fie închis, pentru a putea descrie
o suprafață închisă $ K $! De asemenea, evident, forma diferențială trebuie
să fie de clasă $ \kal{C}^1 $, inclusiv în $ K $, pentru a putea calcula
derivatele parțiale.

\subsection{Forme diferențiale închise}

Fie $ \alpha = Pdx + Qdy $ o 1-formă diferențială de clasă $ \kal{C}^1 $ pe o
vecinătate a $ K = \dr{Int}(\gamma) $. Această formă diferențială se numește
\emph{închisă} dacă are loc:
\[
  \frac{\partial Q}{\partial x} = \frac{\partial P}{\partial y}.
\]

Se poate observa, folosind formula Green-Riemann, că pentru forme diferențiale
închise, integrala curbilinie în lungul oricărui drum este nulă. Această observație
se mai numește \emph{independența de drum a integralei curbilinii} sau
\emph{teorema Poincar\'e}. \index{teorema!Poincar\'e}


%%%%%%%%%%%%%%%%%%%%%%%%%%%%%%%%%%%%%%%%%%%%%%%%%%%%%%%%%%%%%%%%%%%%%%
\section{Exerciții}
1. Să se calculeze următoarele integrale curbilinii de speța întîi:
\begin{enumerate}[(a)]
\item $ \ds\int_\gamma xds $, unde $ \gamma : y = x^2, x \in [0, 2] $;
\item $ \ds\int_\gamma y^5ds $, unde $ \gamma : x = \dfrac{y^4}{4}, y \in [0, 2] $;
\item $ \ds\int_\gamma x^2ds $, unde $ \gamma: x^2 + y^2 = 2, x, y \geq 0 $;
\item $ \ds\int_\gamma y^2ds $, unde $ \gamma: x^2 + y^2 = 4, x \leq 0, y \geq 0 $.
\end{enumerate}

\vspace{1cm}

2. Calculați, direct și aplicînd formula Green-Riemann, integrala curbilinie
$ \ds\int_\gamma \alpha $ în următoarele cazuri:
\begin{enumerate}[(a)]
\item $ \alpha = y^2 dx + x dy $, unde $ \gamma $ este pătratul cu vîrfurile
  $ A(0,0), B(2, 0), C(2, 2), D(0, 2) $;
\item $ \alpha = ydx + x^2 dy $, unde $ \gamma $ este cercul centrat în origine
  și de rază $ 2 $.
\end{enumerate}

\vspace{1cm}

3. Calculați următoarele integrale curbilinii de speța a doua:
\begin{enumerate}[(a)]
\item $ \ds\int_\gamma xdy - ydx $, unde $ \gamma: x^2 + y^2 = 4, y = x\sqrt{3} \geq 0 $
  și $ y = \dfrac{x}{\sqrt{3}} $;
\item $ \ds\int_\gamma (x + y) dx + (x - y) dy $, pe domeniul:
  \[
    \gamma: x^2 + y^2 = 4, y \geq 0;
  \]
\item $ \ds\int_\gamma \dfrac{y}{x + 1} dx + dy $, unde $ \gamma $ este
  triunghiul cu vîrfurile $ A(2, 0), B(0, 0), C(0, 2) $.
\end{enumerate}

\vspace{1cm}

4. Să se calculeze $ \ds\int_\gamma ydx + xdy $ pe un drum de la $ A(2, 1) $
la $ B(1, 3) $.

\vspace{1cm}

5. Să se calculeze circulația cîmpului de vectori $ \vec{V} $ de-a lungul curbei
$ \gamma $, pentru:
\begin{enumerate}[(a)]
\item $ \vec{V} = -(x^2 + y^2) \vec{i} - \vec{x^2 - y^2} \vec{j} $, cu:
  \[
    \gamma: \{ x^2 + y^2 = 4, y < 0 \} \cap \{ x^2 + y^2 - 2x = 0, y \geq 0 \}
  \]
\item $ \vec{V} = x\vec{y} + xy \vec{j} $, unde:
  \[
    \gamma: \{ x^2 + y^2 = 1 \} \cap \{ x + y = 3 \}.
  \]
\end{enumerate}

\vspace{1cm}

6. Să se calculeze masa firului material $ \gamma $, cu ecuațiile parametrice:
\[
  \begin{cases}
    x(t) &= t \\
    y(t) &= t^2/2, t \in [0, 1] \\
    z(t) &= t^3/3
  \end{cases}
\]
iar densitatea $ f(x, y, z) = \sqrt{2y} $.

\vspace{1cm}

7. Fie forma diferențială:
\[
  \alpha = \dfrac{-y}{x^2 + y^2} dx + \dfrac{x}{x^2 + y^2} dy.
\]
Să se calculeze $ \ds\int_\gamma \alpha $, unde $ \gamma $ este cercul centrat
în origine și cu raza 2.

\vspace{1cm}

8. Să se calculeze integrala curbilinie $ \ds\int_\gamma xydx + \dfrac{x^2}{2} dy $,
pe conturul:
\[
  \gamma : \{ x^2 + y^2 = 1, x \leq 0 \leq y \} \cap \{ x + y = -1, x, y < 0 \}.
\]
\emph{Indicație:} Curba $ \gamma $ nu este închisă, deci nu putem aplica
formula Green-Riemann. Considerăm segmentul orientat $ [AB] $, cu
$ A(0, -1), B(0, 1) $, cu care închidem curba. Definim $ C = \gamma \cup [AB] $
și acum putem aplica Green-Riemann pe $ C $. Vom avea, de fapt:
\[
  \int_C \alpha = \int_\gamma \alpha + \int_{[AB]} \alpha,
\]
unde $ \alpha $ este forma diferențială de integrat.

Putem calcula acum integrala pe $ C $ cu Green-Riemann, iar cea pe $ [AB] $ cu
definiția, obținînd în final integrala pe $ \gamma $.

\vspace{1cm}

9. Calculați, folosind integrala curbilinie:
\begin{enumerate}[(a)]
\item lungimea unui cerc de rază 2;
\item lungimea segmentului $ AB $, cu $ A(1, 2) $ și $ B(3, 5) $;
\item lungimea arcului de parabolă $ y = 3x^2 $, cu $ x \in [-2, 2] $;
\item lungimea arcului de hiperbolă $ xy = 1 $, cu $ x \in [1, 2] $.
\end{enumerate}

%%%%%%%%%%%%%%%%%%%%%%%%%%%%%%%%%%%%%%%%%%%%%%%%%%%%%%%%%%%%%%%%%%%%%%
\section{Exerciții suplimentare (ETTI)}

Calculați integralele curbilinii de mai jos:

\textbf{Speța întîi}:
\begin{enumerate}[(a)]
\item $ \ds\int_\gamma ye^{-x} ds $, unde parametrizarea curbei $ \gamma $
  este dată de:
  \[
    \gamma : %
    \begin{cases}
      x(t) &= \ln(1 + t^2) \\
      y(t) &= 2 \arctan t - t
    \end{cases}, \quad t \in [0, 1];
  \]
\item $ \ds\int_\gamma (x^2 + y^2) \ln z ds $, unde parametrizarea curbei
  $ \gamma $ este dată de:
  \[
    \gamma : %
    \begin{cases}
      x(t) &= e^t \cos t \\
      y(t) &= e^t \sin t \\
      z(t) &= e^t
    \end{cases}, \quad t \in [0, 1];
  \]
\item $ \ds\int_\gamma xyz ds $, unde parametrizarea curbei $ \gamma $ este:
  \[
    \gamma : %
    \begin{cases}
      x(t) &= t \\
      & \\
      y(t) &= \dfrac{2}{3}\sqrt{2t^3} \\
      & \\
      z(t) &= \dfrac{1}{2} t^2
    \end{cases}, \quad t \in [0, 1];
  \]
\item $ \ds\int_\gamma x^2y ds $, unde $ \gamma = [AB] \cup [BC] $,
  iar capetele segmentelor sînt $ A(-1, 1), B(2, 1), C(2, 5) $;
\item $ \ds\int_\gamma x^2 ds $, unde $ \gamma $ este dată de:
  \[
    \gamma : x^2 + y^2 = 2, \quad x, y \geq 0;
  \]
\item $ \ds\int_\gamma (x^2 + y^2) ds $, unde $ \gamma $ este sectorul
  de cerc $ x^2 + y^2 = 1 $, parcurs de la $ A(0, -1) $ la $ B(1, 0) $;
\item Calculați lungimea segmentului $ [AB] $, unde $ A(1, 2), B(3, 5) $.
\end{enumerate}

\textbf{Speța a doua}: $ \ds\int_\gamma \omega $ pentru:

\begin{enumerate}[(a)]
\item $ \omega = (x^2 + y^2) dx + (x^2 - y^2) dy $, unde $ \gamma $ este
  dată de:
  \[
    \gamma: %
    \begin{cases}
      x(t) &= \sqrt{t} \\
      y(t) &= \sqrt{t + 1}
    \end{cases}, \quad t \in [1, 4];
  \]
\item $ \omega = \dfrac{1}{y^2 + 1} dx + \dfrac{y}{x^2 + 1} dy $, unde
  $ \gamma $ este dată de:
  \[
    \gamma: %
    \begin{cases}
      x(t) &= t^2 \\
      y(t) &= t
    \end{cases}, \quad t \in [0, 1];
  \]
\item $ \omega = \sqrt{x} dx + xy^2 dy $, unde $ \gamma $ este
  parabola $ y = x^2 $, pentru $ x \in [0, 1] $.
\end{enumerate}

%%% Local Variables:
%%% mode: latex
%%% TeX-master: "../am1-20-21"
%%% End:
